\documentclass[11pt]{article}

\usepackage[margin=1in]{geometry}
\usepackage{amsmath,amssymb}
\usepackage{booktabs}
\usepackage{hyperref}
\usepackage[numbers]{natbib}
\usepackage{xcolor}
\usepackage{microtype}

\newcommand{\pending}{\textcolor{red}{[run pending]}}

\title{\textbf{Hakuli: An Independent Clean-Room Implementation and\\ Reproducibility Study of APOLLO for Memory-Efficient LLM Training}}

\author{
XAEDRO \\
\small\texttt{https://github.com/xaedro/hakuli}
}

\date{October 2026}

\begin{document}

\maketitle

\begin{abstract}
We present \textbf{Hakuli}, an independent, clean-room implementation of the
APOLLO optimizer (Zhu et al., MLSys 2025) released under the Apache License
2.0. Hakuli makes \emph{no claim of algorithmic novelty}: the method follows
the published APOLLO specification, which in turn builds on the low-rank
gradient projection of GaLore (Zhao et al., ICML 2024) and the norm-growth
limiter of Fira (Chen et al., NeurIPS 2025). Our contributions are: (i)~a
self-contained, dependency-free implementation written from the published
description without using reference source code; (ii)~a reproducibility study
of APOLLO's memory and performance claims on consumer hardware; (iii)~a
black-box comparison against the official \texttt{apollo-torch} package under
an identical protocol on a single NVIDIA RTX 3090; and (iv)~measured
full-parameter training of a 7B model at batch size 16 within 24\,GB without
quantization. We also document every deviation from the reference formulation
and provide all scripts, configs and raw results for reproduction.
\end{abstract}

\noindent\textbf{In thirty seconds.} Full-parameter training of 7B language
models is normally out of reach on a 24\,GB consumer GPU because AdamW's
optimizer state alone needs roughly 61\,GB. Hakuli tracks Adam-style moments
in a low-rank random subspace and applies a channel-wise scaling to the full
gradient, reducing optimizer state to a few gigabytes. We rebuilt this method
independently from its public specification, verified it against the official
implementation in a black-box comparison, and measured 7B training at batch
size 16 within 20.5\,GB on an RTX 3090. The software is Apache-2.0.

\section{Introduction}

Memory-efficient optimizers have become a practical requirement for
full-parameter training of large language models (LLMs) on commodity
hardware. AdamW \citep{loshchilov2019adamw} maintains two fp32 moments per
trainable parameter, approximately 8 bytes per parameter, which amounts to
about 61\,GB of optimizer state for a 7B model. Low-rank methods such as
GaLore \citep{zhao2024galore}, Fira \citep{chen2025fira}, APOLLO
\citep{zhu2025apollo} and Q-GaLore \citep{zhang2024qgalore} reduce this
burden by tracking optimizer statistics in a compressed subspace.

APOLLO received an Outstanding Paper Honorable Mention at MLSys 2025 and is
integrated into widely used frameworks. Its official implementation,
\texttt{apollo-torch}, is released under CC-BY-NC, and the Hugging Face
Trainer integration of APOLLO requires that package. For commercial users,
this creates a licensing obstacle independent of the method's technical
merit.

This paper presents Hakuli, an independent, clean-room implementation of the
APOLLO method under Apache-2.0. We emphasize what this work is and is not:

\begin{itemize}
  \item \textbf{Not} a new optimizer. The update rule follows APOLLO's
  published specification. We explicitly reject any framing of Hakuli as an
  algorithmic invention.
  \item \textbf{An} independent implementation: single file, PyTorch-only, no
  reference source code consulted, numerically validated against a
  specification-derived reference before any official code was executed.
  \item \textbf{A} reproducibility study: we test which of APOLLO's reported
  memory and performance claims can be independently reproduced on a single
  24\,GB consumer GPU, and which cannot be tested with our hardware.
  \item \textbf{A} black-box comparison: identical protocol, same GPU, official
  \texttt{apollo-torch} v1.0.3 versus Hakuli.
  \item \textbf{A} documented set of deviations from the reference formulation
  (Section~\ref{sec:deviations}) and a full reproducibility package.
\end{itemize}

\section{Background: APOLLO}
\label{sec:background}

We summarize the method as published; notation follows the APOLLO paper
\citep{zhu2025apollo}.

For a weight matrix with gradient $G_t \in \mathbb{R}^{m \times n}$ and a
random projection $P_t \in \mathbb{R}^{m \times r}$ resampled every
\texttt{update\_proj\_gap} steps, APOLLO computes a projected gradient
$R_t = P_t^\top G_t$ and tracks first and second moments in the projected
space,
\begin{align}
  M_t &= \beta_1 M_{t-1} + (1-\beta_1) R_t, \\
  V_t &= \beta_2 V_{t-1} + (1-\beta_2) R_t^2 .
\end{align}
The Adam-normalized direction in the subspace is
$S_t = M_t / (\sqrt{V_t} + \epsilon)$, and a channel-wise scaling factor is
obtained from the ratio of norms,
\begin{equation}
  s_t = \frac{\lVert S_t \rVert_{\mathrm{dim}}}{\lVert R_t \rVert_{\mathrm{dim}}},
  \qquad
  U_t = \sqrt{\lambda}\, s_t \odot G_t,
\end{equation}
where the norms are taken along the output-channel dimension and $\lambda$ is
a heuristic scale. The update is applied to the \emph{full} gradient with a
scale-corrected step size,
\begin{equation}
  \eta_t = \eta \, \frac{\sqrt{1-\beta_2^t}}{1-\beta_1^t},
  \qquad
  W_t = W_{t-1} - \eta_t U_t .
\end{equation}
A tensor-wise variant (APOLLO-Mini) replaces the channel vector $s_t$ by a
scalar and uses rank $r=1$. The norm-growth limiter (NL), adopted from Fira
\citep{chen2025fira}, caps the relative growth of the update norm between
consecutive steps and stabilizes training.

The lineage is direct: GaLore introduced low-rank gradient projection with a
periodic SVD; Fira introduced norm-based scaling with a norm-growth limiter;
APOLLO replaced the SVD with a pure random projection and moved the update
back to the full gradient. Hakuli implements this final formulation.

\section{Clean-Room Methodology}
\label{sec:cleanroom}

\textbf{Implementation provenance.} Hakuli was implemented independently from
the published APOLLO algorithmic description (arXiv:2412.05270) and its
public documentation. No source code from the official APOLLO implementation
was used in the implementation of Hakuli.

Concretely:

\begin{itemize}
  \item The implementation was written from the mathematical specification
  and validated numerically against an internal, specification-derived
  reference implementation before any official code was executed.
  \item The first execution of the official implementation
  (\texttt{apollo-torch} v1.0.3) took place on 2026-10-09, as a black box,
  for the comparison reported in Section~\ref{sec:reproduction} --- after the
  Hakuli implementation was complete (initial public commit
  \texttt{a8ef673}).
  \item The official package is used exclusively as an external dependency of
  the benchmark harness for black-box runs. It is not linked, copied, or
  redistributed by Hakuli.
  \item Internal records (dates, commit history, and the development log of
  the specification-derived validation) are retained by the authors.
\end{itemize}

This separation of roles --- implementation from the specification first,
black-box validation against the reference second --- is what makes the
clean-room claim verifiable in practice.

\section{Hakuli Implementation}
\label{sec:implementation}

Hakuli is a single-class \texttt{torch.optim.Optimizer} subclass with no
dependencies beyond PyTorch ($\geq$2.1). Design choices:

\begin{itemize}
  \item \textbf{Random projection with periodic resampling}, no SVD at any
  step; the projection is stored in bfloat16.
  \item \textbf{Per-parameter instances with per-layer weight updates.} One
  optimizer instance per low-rank parameter, stepped from
  \texttt{register\_post\_accumulate\_grad\_hook} so gradients are released
  eagerly. This pattern was introduced in the GaLore codebase and adopted by
  APOLLO and Q-GaLore; we follow it and cite it as such.
  \item \textbf{Adam-style fallback.} Parameters whose smallest dimension does
  not exceed the rank, and all 1-D parameters, are updated with a
  bias-free Adam-style direction. A shared AdamW covers these parameters.
  \item \textbf{Decoupled weight decay}, standard \texttt{state\_dict}
  checkpointing, and closure support.
  \item \textbf{Optional loss-memory reduction.} In our training runs we pair
  Hakuli with Cut Cross-Entropy \citep{wijmans2025cce}, which computes the
  cross-entropy without materializing the full logit matrix.
\end{itemize}

Per projected parameter of shape $m \times n$, optimizer state is
$2 r \max(m,n)$ bytes (bf16 projection) plus $8 r \min(m,n)$ bytes (two fp32
moments), versus $8mn$ bytes for AdamW fp32 moments. For a 7B transformer at
$r=256$ this is approximately 2.4\,GB in total.

\section{Experimental Setup}
\label{sec:setup}

\textbf{Hardware and software.} Single NVIDIA RTX 3090 (24\,GB), PyTorch
2.5.1+cu124, Transformers 4.57.6, Datasets 5.0.1, SDPA attention, gradient
checkpointing enabled, bf16 weights and gradients.

\textbf{Model and data.} LLaMA-7B (\texttt{huggyllama/llama-7b}), trained on
streaming C4 \citep{raffel2020t5} with packed sequences. We start from the pretrained
weights (continual pretraining); the reported loss is therefore not
comparable to from-scratch runs and serves only as a sanity check that
training is healthy.

\textbf{Protocol (controlled comparison).} Batch size 16, sequence length
256, 300 optimizer steps, learning rate $5\times10^{-3}$, rank 512,
\texttt{update\_proj\_gap} 500, 10 warmup steps. The official
\texttt{apollo-torch} (v1.0.3) and Hakuli are run with this identical
configuration in the same harness.

\textbf{Measurement.} Peak VRAM is \texttt{torch.cuda.max\_memory\_allocated}
after resetting statistics at the end of warmup, so one-time initialization
allocations are excluded. Throughput is steady-state tokens/second measured
over steps after warmup. Raw JSON results, configs and the harness are
included in the repository (\texttt{benchmarks/}).

\section{Reproduction of APOLLO}
\label{sec:reproduction}

We attempted to reproduce the reported claims that are testable on a single
24\,GB GPU. Table~\ref{tab:claims} summarizes the status.

\begin{table}[t]
\centering
\small
\begin{tabular}{p{0.40\textwidth}p{0.20\textwidth}p{0.34\textwidth}}
\toprule
APOLLO claim & Status here & Evidence / reason \\
\midrule
7B pre-training within 11\,GB (A100, batch 1) & partially addressed &
Different batch size and GPU; we measure at batch 16 on 24\,GB (Table~\ref{tab:main}). \\
Up to $3\times$ throughput vs.\ Adam on $4\times$A100-80GB & not testable &
Hardware unavailable. \\
LLaMA-13B with naive DDP on A100-80GB & not testable & Hardware unavailable. \\
7B under 12\,GB combined with quantization & not addressed &
Hakuli v1.0 implements no quantization. \\
Structured (channel-wise) scaling matches AdamW-level quality & consistent &
Healthy loss trajectories in our runs; no quality regression observed. \\
\bottomrule
\end{tabular}
\caption{Status of APOLLO's reported claims under our hardware constraints.}
\label{tab:claims}
\end{table}

\textbf{Black-box comparison.} Under the protocol of Section~\ref{sec:setup},
the official \texttt{apollo-torch} v1.0.3 peaked at 19.44\,GB and reached
1058.1\,tokens/s (final loss 6.73); Hakuli peaked at
\textbf{20.51\,GB} and reached \textbf{1065.1}\,tokens/s (final loss 2.36).
Both implementations thus fit 7B full-parameter training at batch size 16
within 24\,GB on this GPU. The final losses are not directly comparable: the
two implementations apply their scale hyperparameters under different
conventions, and each was run with its own documented default heuristic;
memory and throughput are the controlled quantities. Remaining differences
between the two rows reflect implementation details listed in
Section~\ref{sec:deviations} and run-to-run variation, not algorithmic
divergence.

\section{Systems Performance}
\label{sec:systems}

\begin{table}[t]
\centering
\small
\begin{tabular}{llrrrll}
\toprule
Method & GPU & Batch & Rank & State precision & Peak VRAM & Source \\
\midrule
AdamW & 24\,GB class & 16 & --- & fp32 & does not fit & \citep{zhao2024galore} \\
GaLore-8bit & RTX 4090 & 16 & 1024 & int8 & 22.8\,GB & GaLore repo \\
Q-GaLore & A6000 & 16 & 1024 & int4 & 15.26\,GB & Q-GaLore repo \\
APOLLO (official) & RTX 3090 & 16 & 512 & bf16/fp32 & 19.44\,GB & this work \\
\textbf{Hakuli} & RTX 3090 & 16 & 512 & bf16/fp32 & \textbf{20.51\,GB} & this work \\
\bottomrule
\end{tabular}
\caption{Peak training memory for LLaMA-7B on 24\,GB-class GPUs. Published
anchors were measured by their authors on different hardware, ranks and
state precisions; they are included for context, not as controlled
comparisons. The controlled comparison is the last two rows (same protocol,
same GPU).}
\label{tab:main}
\end{table}

Table~\ref{tab:main} summarizes peak memory. The controlled rows show that
both the official implementation and Hakuli enable LLaMA-7B full-parameter
training at batch size 16, sequence length 256, within 24\,GB, without
quantization. Hakuli's measured throughput of 1065\,tokens/s corresponds to
about 1.6\,s per optimizer step at 4096 tokens per step.

\textbf{Long-context fine-tuning.} In a separate internal run on the same
GPU, a 7B model was full-parameter fine-tuned at sequence length 3072, batch
size 4, rank 256, using Hakuli with Cut Cross-Entropy, peaking at 23.3\,GB.
This demonstrates that the memory profile holds at long context, where
activation memory dominates.

\section{Comparison with Related Methods}
\label{sec:comparison}

Beyond the controlled comparison, we relate Hakuli to the broader family of
memory-efficient optimizers:

\begin{itemize}
  \item \textbf{AdamW}: best quality baseline, but optimizer state alone
  ($\approx$61\,GB for 7B) exceeds consumer GPUs.
  \item \textbf{GaLore}: SVD-based projection; periodic decomposition stalls
  and degraded stability at small batch sizes are well documented in later
  work \citep{zhu2025apollo}.
  \item \textbf{Q-GaLore}: adds INT4 projection and quantization, reaching
  15.26\,GB for 7B at batch 16, at the cost of a quantized state.
  \item \textbf{SCALE} \citep{glentis2026scale}: a 2026 minimalist design
  (column-wise gradient normalization plus last-layer momentum) reporting
  better perplexity and memory than APOLLO for 7B. We cite it as the current
  state of the art in this line and as future comparison work; it was not
  benchmarked here.
\end{itemize}

\section{Ablations}
\label{sec:ablations}

This version does not include a systematic ablation study. We report the
observations we have and mark the rest as planned work:

\begin{itemize}
  \item \textbf{Rank.} Rank 512 at batch size 16 fits 24\,GB (20.5\,GB
  measured); rank 256 was used in our long-context run (23.3\,GB at
  sequence length 3072, where activations dominate). A rank sweep at fixed
  sequence length is planned.
  \item \textbf{Projection interval.} \texttt{update\_proj\_gap}=500 was used
  throughout; no sensitivity analysis yet.
  \item \textbf{Scale.} The heuristic $\lambda \approx
  \mathrm{hidden}/\mathrm{rank}$ was used in fine-tuning runs; no sweep yet.
  \item \textbf{Precision.} bfloat16 projections were used in all runs; a
  comparison against fp32 projections is planned.
\end{itemize}

\section{Limitations and Deviations from APOLLO}
\label{sec:deviations}

\textbf{Limitations.} All controlled results come from a single GPU (RTX
3090) and a single model size (7B); multi-GPU and FSDP behavior of Hakuli is
untested. Loss values are reported for protocol sanity, not for quality
claims. No quantization is implemented.

\textbf{Deviations from the reference formulation.} Hakuli differs from the
official implementation in the following, all documented in the code:

\begin{enumerate}
  \item The projection matrix is stored in bfloat16.
  \item The \texttt{scale} parameter is applied under a square root
  ($\sqrt{\lambda}$), matching our validated internal convention.
  \item The norm-growth limiter is expressed as a relative cap with a
  restart when the previous update norm is exactly zero; the official
  implementation applies the limiter before or after scaling depending on a
  configuration flag.
  \item The Adam fallback applies when $\min(m,n) \leq r$ or the parameter is
  not 2-D; the reference partitions parameters differently.
  \item Per-parameter optimizer instances plus a shared AdamW for non-2-D
  parameters; the reference supports both single-optimizer and per-layer
  configurations.
  \item No quantization or SVD projection options are implemented.
\end{enumerate}

\section{Reproducibility Statement}
\label{sec:reproducibility}

The repository contains the optimizer, a test suite (unit and integration
tests, including a guard that no SVD routine is ever called), the benchmark
harness, all JSON configs, and the raw result files reported here.
Reproduction:

\begin{verbatim}
git clone https://github.com/xaedro/hakuli
pip install -r benchmarks/requirements.txt
python benchmarks/run_benchmark.py \
  --config benchmarks/configs/memory_llama7b_hakuli.json
python benchmarks/table.py
\end{verbatim}

The official comparison additionally requires \texttt{pip install
apollo-torch} (v1.0.3) and the config
\texttt{memory\_llama7b\_apollo.json}. Environment: Docker container based on
NVIDIA NGC PyTorch 24.10, Python 3.11. Initial public commit:
\texttt{a8ef673}.

\section{License and Provenance}
\label{sec:license}

Hakuli is released under the Apache License 2.0, including a patent grant
from its contributors over their own contributions. This grants no rights to
third-party patents. The official APOLLO implementation is released under
CC-BY-NC (with portions from GaLore under Apache-2.0); GaLore and Fira are
Apache-2.0. Hakuli is not affiliated with, or endorsed by, the APOLLO
authors. Academic attributions are recorded in the repository \texttt{NOTICE}
file. We thank the APOLLO, GaLore and Fira authors for releasing their
implementations and documentation, which made this study possible.

\bibliographystyle{plainnat}
\bibliography{references}

\end{document}
